% Source: Solutions/Sahil-PS1-Q6-10.pdf, p. 2 (typed)
\submittedby{Sahil}

Let $\mathcal{B}$ be a basis for a topology on $X$, and let $\T_{\mathcal{B}}$ denote the
topology generated by $\mathcal{B}$. By the definition of the topology generated by a basis,
\[ \T_{\mathcal{B}}=\Bigl\{\bigcup_{\alpha\in A}B_\alpha : A \text{ is an index set and }
   B_\alpha\in\mathcal{B}\Bigr\}. \]
Let $\mathfrak{F}$ be the collection of all topologies on $X$ that contain $\mathcal{B}$. We prove
\[ \boxed{\T_{\mathcal{B}}=\bigcap_{\T\in\mathfrak{F}}\T.} \]

First, let $U\in\T_{\mathcal{B}}$. Then $U$ is a union of elements of $\mathcal{B}$. If
$\T\in\mathfrak{F}$, then $\mathcal{B}\subseteq\T$. Since every topology is closed under
arbitrary unions, it follows that $U\in\T$. This holds for every $\T\in\mathfrak{F}$, and hence
\[ \T_{\mathcal{B}}\subseteq\bigcap_{\T\in\mathfrak{F}}\T. \]

For the reverse inclusion, note that $\T_{\mathcal{B}}$ is itself a topology on $X$ containing
$\mathcal{B}$. Consequently, $\T_{\mathcal{B}}\in\mathfrak{F}$. The intersection of all members
of $\mathfrak{F}$ is therefore contained in this particular member:
\[ \bigcap_{\T\in\mathfrak{F}}\T\subseteq\T_{\mathcal{B}}. \]
Combining the two inclusions proves the required equality.

\textbf{Interpretation:} The topology generated by a basis is the smallest topology on $X$ that
contains every basis element.
